% section: はじめに

  \section{はじめに}
    \subsection{本書について}
      本書は漸化式の解法を体系的にまとめ,それぞれについて解説するものである.
      解説は解法の発想を明らかにし,流れに沿った論理の展開と行間の少ない設計を目指した.
      それらの確認と定着のため,各解法は,いくつかの例題を通じて解説している.
      また,その基礎として最低限の計算上のテクニックの解説も添えた.
      例題はその種の漸化式において解けなければいけないレベルのものを3題前後用意し,例題に対して,方針,解答,解法を付した.
      それぞれ方針は解法の手がかりを,解答は一般項の値を,解法は結果に至る過程を扱っている.

      本書を読む際は,机,紙,ペン,それときちんと一行一行を読む時間を用意してからページをめくってほしい.
      教科書を読むいうのは案外難しいもので,議論を丁寧に追うのは相応の苦労を伴う.
      しかしそれを経て得た知識は,演習によって強化され今後の糧となる.
      そのために,一行一行をよく読みここに書かれていることを実際に自身で試してほしい.

      本書を読むことで,与えられた漸化式の一般項を得るにあたって必要な手立ての検討を行えるようになる.
      きちんと読みこみ,十分な演習を行えば大学受験程度の漸化式においては得意分野と言えるぐらいにはなるだろう.     

    \subsection{本書の内容と対象読者}
      本書は漸化式を一通り解けることを目標としている.
      それに際して基本的な漸化式の解法から話をする.
      最終的には難関国公立大学に出題されるような漸化式や,数学的に興味深い性質を持つ漸化式に至るまで,初歩から発展まで幅広く扱っている.
      
      対象読者は数学が好きだが,数列の分野については演習が不十分な高校生である.
      従って大部分は高校数学の初歩的な知識と最低限の計算力を前提に進める.

    \subsection{本書でのルール}
      特に断りが無い限りは,\,$d,i,k,l,m,n,$は1以上の整数,\,$p,q,r,s,t,u,v,x,y,z,\alpha,\beta,\gamma,\delta,\lambda$は実数として扱う.
      また,\,$\therefore\,,\because\,,\Rightarrow\,,\Leftrightarrow$は左から順に「従って」「なぜならば」「左側が成り立つならば右側が成り立つ」「同値である」を意味する.
      加えて,\,$\cdots$や$\vdots$は途中を省略する時に用いる.

